\SourceLeaf{96}%

\noindent This is one of the series of theorems that \textsc{Fourier} discovered from 1807 onward and published in detail in 1822 in the »Théorie analytique de la chaleur«. The summation formula given by \textsc{Gauss} to prove the theorem yields, for $\operatorname{lim.}\,\omega=0$, \textsc{Fourier}'s integral formula directly. Observe that the equation $\mathop{\mathrm{lim.}}\limits_{x=\pm\infty}\varphi(x)=0$ must hold here. The function that \textsc{Gauss} uses in probability theory, $\varphi(x)$, does satisfy this condition (cf. »Theoria combinationis observationum etc.« Art. 4, Collected Works IV p. 5).

\Needspace{17\baselineskip}
As for \textsc{Gauss}'s summation formula itself, splitting the integral on the right first puts it in the form
\[
\sum^{+\infty}_{k=-\infty}\varphi\!\left(t+\frac{2\pi k}{\omega}\right) = \frac{\omega}{2\pi}\sum^{+\infty}_{l=-\infty}\left(e^{-ilt\omega}\mathord{.}\sum^{+\infty}_{k=-\infty}\int^{\frac{2(k+1)\pi}{\omega}}_{\frac{2k\pi}{\omega}}\varphi(\theta)\mathord{.} e^{il\theta\omega}\,d\theta\right).
\]
\Needspace{15\baselineskip}
\noindent Introducing a new variable of integration $x$ by $\theta\omega = x+2k\pi$ and reversing the order of the sums on the right gives the following new form of the summation formula.
\[
\sum^{\infty}_{k=-\infty}\varphi\!\left(t+\frac{2\pi k}{\omega}\right) = \sum^{+\infty}_{k=-\infty}\sum^{+\infty}_{l=-\infty}\left(e^{-ilt\omega}\mathord{.}\frac{1}{2\pi}\int^{2\pi}_{0}\varphi\!\left(\frac{x+2k\pi}{\omega}\right)e^{ilx}\,dx\right).
\]
\Needspace{12\baselineskip}
\noindent The validity of this formula follows directly from the \textsc{Fourier} series.
\[
\varphi\!\left(t+\frac{2\pi k}{\omega}\right) = \sum^{+\infty}_{l=-\infty}\left(e^{-ilt\omega}\mathord{.}\frac{1}{2\pi}\int^{2\pi}_{0}\varphi\!\left(\frac{x+2k\pi}{\omega}\right)e^{ilx}\,dx\right).
\]

\begin{flushright}
\textsc{Fricke}.
\end{flushright}
